% !PRINTED 429
\weberpage{429}
%<<CONT
ist aus $\varphi_1\varphi_2$ komponiert. Die Form $\varphi$ gehört in eine Klasse $k$, die aus $k_1,k_2$ \gesperrt{komponiert} heißt.

Nehmen wir $\varphi=(a,b,c)$ so an, daß $a$ relativ prim zu $m$ ist, so ist auch $b$ relativ prim zu $a$ (wegen $m=4ac-b^2$). Ist $a=a_1a_2$, so sind auch $(a_1,b,ca_2)$, $(a_2,b,ca_1)$ einhellige primitive Formen, und man erhält die Komposition
\[
(a,b,c) = (a_1,b,ca_2),\ (a_2,b,ca), % [sic] source omits the expected subscript 1 on the final a.
\]
und durch wiederholte Anwendung dieses Satzes gelangt man zu dem Resultat:

I. \gesperrt{Die Form $\varphi=(a,b,c)$ läßt sich, wenn $a$ relativ prim zu $m$ ist, aus solchen Formen zusammensetzen, deren erste Koeffizienten Primzahlen, nämlich die Primfaktoren von $a$ sind.}

Eine dieser Komponenten, etwa
\[
(p,b,acp^{-1}),
\]
läßt sich durch eine äquivalente Form $(p,b',c'p^g)$ ersetzen, deren dritter Koeffizient $c'p^g$ durch eine beliebig hohe Potenz von $p$ teilbar ist.

Um dies zu beweisen, brauchen wir für die Äquivalenz das Zeichen $\sim$, und haben
\[
(p,b,cp^{g-1}) \sim (p,b+2\lambda p^g,c'p^g),
\]
worin $c'$ wegen der Gleichheit der Diskriminanten aus
\[
pc' = c+\lambda b+\lambda^2 p^g
\]
zu bestimmen ist. Da nun $b$ durch $p$ nicht teilbar ist, so kann $\lambda$ aus der Kongruenz $c+\lambda b \equiv 0\ (\mathrm{mod}\ p)$ bestimmt werden, und $c'$ ergibt sich als ganze Zahl. Damit ist mit Rücksicht auf (I.) bewiesen:

II. \gesperrt{Man kann in jeder Klasse $k$ von primitiven Formen der Diskriminante $D$ einen Repräsentanten $\varphi$ finden, der sich aus Formen}
\[
P=(p,b,p^gc),
\]
\gesperrt{worin $p$ eine in $D$ nicht aufgehende Primzahl, und $g$ ein beliebiger Exponent ist, zusammensetzen läßt.}

Die Form $P$ läßt sich leicht mit sich selbst zusammensetzen, denn es ist im Sinne der Kompositionen, wenn $g>\nu$ ist:
\[
(p^\nu,b,p^{g-\nu}c)(p,b,p^gc) = (p^{\nu+1},b,p^{g-\nu-1}c),
\]
und folglich durch den Schluß von $\nu-1$ auf $\nu$:
\weq{2}{P^\nu = (p^\nu,b,p^{g-\nu}c).}
